\scanpage{11}
\mspage{32}
\figure{\begin{tikzpicture}\draw[->](0,2.4)--(0,0);\end{tikzpicture}}
\deleted{
L'assertion (iii) se traite d'une façon analogue. Puisque les foncteurs $g^*(R^qf_*)$ sont effaçables pour $q>0$ dans la catégorie des faisceaux de torsion (8.1), la bijectivité de $\varphi^n$ pour $F$ implique l'effaçabilité de $(R^nf'_*)g'^*$ pour $F$. Inversement, supposons que $\varphi^{n-1}$ soit bijectif pour $(g,f,.)$ et que $(R^nf'_*)g'^*$ soit effaçable pour le faisceau $F$. Soit d'abord $F\longrightarrow G$ un effacement de $g^*(R^nf_*)$ avec $G$ de torsion et soit $C=G/F$. Alors on peut compléter le diagramme
\[
(*)
\]
\diagram{\begin{tikzcd}[column sep=small]
g^*(R^{n-1}f_*\deleted{\unsure{)}}\inserted{(}G\inserted{)})\arrow[r]\arrow[d,"\unsure{\sim|2}"']&g^*(R^{n-1}f_*\deleted{\unsure{)}}\inserted{(}C\inserted{)})\arrow[r]\arrow[d,"\unsure{\sim|2}"']&g^*(R^nf_*\deleted{\unsure{)}}\inserted{(}F\inserted{)})\arrow[r]\arrow[d,"\varphi^n"]&{}\\
\deleted{\unsure{(}}R^{n-1}f'_*\deleted{\unsure{)}}\inserted{(}g'^*G\inserted{)}\arrow[r]&\deleted{\unsure{(}}R^{n-1}f'_*\deleted{\unsure{)}}\inserted{(}g'^*C\inserted{)}\arrow[r]&\deleted{\unsure{(}}R^nf'_*\deleted{\unsure{)}}\inserted{(}g'^*F\inserted{)}\arrow[r]&{}
\end{tikzcd}}
avec un $0$ au bout de la première ligne, ce qui implique
}
\figure{\begin{tikzpicture}\draw[->](0,2.8)--(0,0);\end{tikzpicture}}
\mspage{10}
